\makeatletter \def\input@path{{../../}} \makeatother
\documentclass[11pt]{article}

\usepackage[utf8]{inputenc}
\usepackage{graphicx}
\usepackage[dvipsnames]{xcolor}
\usepackage{color}
\usepackage[margin=1in]{geometry}
\usepackage{hyperref}
\usepackage{tikz}
\usepackage{amsmath}
\usepackage{commath}
\usepackage{relsize}
\usepackage{centernot}
\usepackage{amssymb}
\usepackage{amsfonts}
\usepackage{graphicx}
\usepackage{textcomp}
\usepackage{gensymb}
\usepackage{enumitem}
\usepackage{siunitx}
\usepackage{fancyhdr}
\usepackage{floatrow}
\usepackage{wrapfig}
\usepackage[labelfont=bf]{caption}
\usepackage{mathtools}
\usepackage{pgfplots}
\usepackage{multicol}
\usepackage[misc]{ifsym}
\usepackage{tabularx}
\usepackage{empheq}
\usepackage{tkz-tab}
\usepackage{xpatch}
\usepackage{caption}
\usepackage{subcaption}
\usepackage{tabularray}
\usepackage{esint}
\usepackage{amsthm}
\usepackage{thmtools}
\RequirePackage[most]{tcolorbox}

\swapnumbers

\makeatletter
\declaretheoremstyle[
  headfont= \sffamily\bfseries\color{cyan!50!black},
  headpunct={\sffamily\bfseries.},
  postheadspace=0.5em,
  notefont=\sffamily\bfseries,
  headformat=\NAME\NUMBER\let\thmt@space\@empty\NOTE,
  bodyfont=\mdseries\slshape,
  spaceabove=12pt,spacebelow=12pt,
]{thmstyle}
\makeatother

\theoremstyle{thmstyle}
\declaretheorem[name=Definici\'on]{theorem}
\newtheorem{definition}[theorem]{Definici\'on}

\definecolor{lightcyan}{rgb}{0.88, 1.0, 1.0}
\colorlet{mythmback}{lightcyan!40!white}

\tcolorboxenvironment{theorem}{
  colback=mythmback,coltitle=blue,colframe=mythmback,
  left=0pt,right=0pt,
  top=0pt,bottom=0pt,
  before skip=10pt, after skip=10pt,
}

\xpatchcmd{\tkzTabLine}{$0$}{$\bullet$}{}{}
\tikzset{t style/.style={style=solid}}
\pgfplotsset{compat=newest}

\newtheorem{manualtheoreminner}{}
\newenvironment{manualtheorem}[1]{%
  \renewcommand\themanualtheoreminner{#1}%
  \manualtheoreminner
}{\endmanualtheoreminner}

\renewcommand{\pd}[2]{\frac{\partial #1}{ \partial #2}}
\newcommand{\td}[2]{\frac{\mathrm{d} #1}{ \mathrm{d} #2}}
\newcommand{\pdd}[2]{\frac{\partial^2 #1}{ \partial #2 ^2}}
\newcommand{\pddc}[3]{\frac{\partial^2 #1}{ \partial #2 \partial #3}}
\newcommand{\solution}[0]{\large{\textbf{Soluci\'on}}\normalsize}

\def\sca   #1{\mbox{\rm{#1}}{}}
\def\mat   #1{\mbox{\boldmath $\mathsf #1$}}
\def\vec   #1{\mbox{\boldmath $#1$}{}}
\def\ten   #1{\mbox{\boldmath $#1$}{}}
\def\ltr   #1{\mbox{\sf{#1}}}
\def\bltr  #1{\mbox{\sffamily{\bfseries{{#1}}}}}

\DeclareMathOperator{\dive}{div}
\DeclareMathOperator{\trace}{trace}
\DeclareMathOperator{\tr}{tr}
\DeclareMathOperator{\symm}{symm}
\DeclareMathOperator{\sk}{skew}
\DeclareMathOperator{\grad}{grad}
\DeclareMathOperator{\Grad}{Grad}
\DeclareMathOperator{\curl}{curl}
\DeclareMathOperator{\Curl}{Curl}
\DeclareMathOperator*{\argmin}{\arg\!\min}

\def\dx{\mbox{\(\,\mathrm{d}x\)}}
\def\ds{\mbox{\(\,\mathrm{d}s\)}}
\def\dS{\mbox{\(\,\mathrm{d}S\)}}
\def\dV{\mbox{\(\,\mathrm{d}V\)}}
\def\dv{\mbox{\(\,\mathrm{d}v\)}}
\def\R{\mathbb R}
\def\C{\mathbb C}
\def\N{\mathbb N}
\def\Q{\mathbb Q}
\def\Z{\mathbb Z}
\def\D{\mathcal D}

\usepackage{mdframed}
\mdfdefinestyle{MyFrame}{%
    linecolor=black,
    linewidth=1pt,
    outerlinewidth=1pt,
    roundcorner=20pt,
    innertopmargin=\baselineskip,
    innerbottommargin=\baselineskip,
    innerrightmargin=20pt,
    innerleftmargin=20pt,
    splittopskip=2\baselineskip,
    backgroundcolor=gray!50!white}

\setlength\textwidth{7in}
\setlength\parindent{0pt}
\oddsidemargin=-0.5cm
\pagestyle{fancy}
\newcommand*{\vertbar}{\rule[-1ex]{0.5pt}{2.5ex}}
\setlength{\headheight}{13.59999pt}

\renewcommand{\headrulewidth}{0.1pt}
\renewcommand{\footrulewidth}{0.1pt}
\renewcommand{\figurename}{Figura}
\renewcommand\refname{Referencias}
\renewcommand{\thesection}{\arabic{section}.}
\renewcommand{\thesubsection}{\thesection\arabic{subsection}.}
\renewcommand{\labelitemi}{{\raisebox{.3\height}{\scalebox{0.6}{$\blacklozenge$}}}}
\renewcommand*\contentsname{\'Indice}

\newcommand{\p}{\vspace{10pt}}

\AtBeginDocument{\vspace*{-3.2\baselineskip}}
\textheight=665pt

\lhead{\scshape Pontificia Universidad Cat\'olica de Chile}
\rhead{\scshape IMC UC}
\cfoot{\thepage}

\usepackage{footnote}
\usepackage{etoolbox}
\newtoggle{indefintion}
\togglefalse{indefintion}
\pretocmd{\footnote}{\iftoggle{indefintion}{\stepcounter{footnote}}{\relax}}{}{}

\begin{document}
    \begin{flushleft}
        \small
        \vspace{1pt}
        \textbf{IMT3130} \hfill
        \textbf{16/04/2026}
    \end{flushleft}
    \begin{center}
        \LARGE\textbf{Ayudant\'ia 5}\\
        \vspace{3pt}
        \normalsize\textbf{Lifting, Galerkin y estimaciones abstractas}\\
        \normalsize Ayudante: Basti\'an Herrera \Letter{\hspace{1pt}: \texttt{biherrera@uc.cl}}
        \rule{\textwidth}{0.05mm}
    \end{center}

\section*{Problemas}

\begin{enumerate}[label=\textbf{\arabic*.}, leftmargin=*]

    \item \textbf{Condiciones mixtas y lifting sobre $\Gamma_D$.} Considere el problema
    \[
    \begin{aligned}
        -\Delta u + cu &= f &&\text{en }\Omega,\\
        \gamma_0u &=g &&\text{en }\Gamma_D,\\
        \partial_n u+\beta\gamma_0u &= r &&\text{en }\Gamma_R,
    \end{aligned}
    \]
    en un dominio Lipschitz acotado con $\partial\Omega=\overline{\Gamma_D}\cup\overline{\Gamma_R}$ y $\Gamma_D\cap\Gamma_R=\varnothing$. Defina
    \[
        V=\{v\in H^1(\Omega):\gamma_0v=0\text{ en }\Gamma_D\},
        \qquad
        V_g=\{v\in H^1(\Omega):\gamma_0v=g\text{ en }\Gamma_D\}.
    \]
    Suponga $f\in L^2(\Omega)$, $r\in H^{-1/2}(\Gamma_R)$, $g\in H^{1/2}(\Gamma_D)$ y $c,\beta\in L^\infty$ con $c,\beta\ge 0$.
    \begin{enumerate}[label=(\alph*)]
        \item Considere primero el caso homog\'eneo $g=0$. Derive la formulaci\'on d\'ebil en el espacio $V$, identificando la forma bilineal y la forma lineal. Suponiendo $|\Gamma_D|>0$, muestre que el problema es bien puesto por Lax--Milgram.
        \item Vuelva ahora al caso general $g\in H^{1/2}(\Gamma_D)$. Tome un lifting $G\in H^1(\Omega)$ tal que $\gamma_0G=g$ en $\Gamma_D$, escriba $u=u_0+G$ con $u_0\in V$, y derive el problema equivalente satisfecho por $u_0$. Identifique con cuidado la nueva funcional lineal y explique por qu\'e la existencia y unicidad de $u_0$ implican existencia y unicidad de $u$ en $V_g$.
    \end{enumerate}

    \item \textbf{Ortogonalidad de Galerkin como proyecci\'on.} Sea $V$ un espacio de Hilbert, $a:V\times V\to\R$ una forma bilineal sim\'etrica, continua y coerciva, y $\ell\in V'$. Sea $u\in V$ la soluci\'on de
    \[
        a(u,v)=\ell(v)
        \qquad\forall v\in V.
    \]
    Sea $V_h\subset V$ un subespacio finito-dimensional conforme y suponga que $u_h\in V_h$ satisface
    \[
        a(u_h,v_h)=\ell(v_h)
        \qquad\forall v_h\in V_h.
    \]
    Usaremos la ortogonalidad de Galerkin, ya demostrada en clases:
    \[
        a(u-u_h,v_h)=0
        \qquad\forall v_h\in V_h.
    \]
    \begin{enumerate}[label=(\alph*)]
        \item Defina $\|v\|_a\coloneqq a(v,v)^{1/2}$ y pruebe que es una norma equivalente a la norma de $V$.
        \item Use la ortogonalidad para concluir que $u_h$ es la mejor aproximaci\'on de $u$ en $V_h$ respecto de $\|\cdot\|_a$, es decir,
        \[
            \|u-u_h\|_a=\inf_{v_h\in V_h}\|u-v_h\|_a.
        \]
        \item Si $V_h=\operatorname{span}\{\varphi_1,\dots,\varphi_N\}$ y $u_h=\sum_{j=1}^N U_j\varphi_j$, escriba el sistema lineal que satisfacen los coeficientes $U_j$.
    \end{enumerate}

\item \textbf{Recuperar la formulaci\'on fuerte desde una formulaci\'on d\'ebil con coeficiente anisotr\'opico y condiciones mixtas.}
Sea $\Omega\subset\R^d$ un dominio Lipschitz acotado, con descomposici\'on
\[
\partial\Omega=\overline{\Gamma_D}\cup\overline{\Gamma_R},
\qquad
\Gamma_D\cap\Gamma_R=\varnothing,
\qquad
|\Gamma_D|>0.
\]
Sea $\mat B\in W^{1,\infty}(\Omega;\R^{d\times d})$ una matriz sim\'etrica, y defina
\[
\mat A_\varepsilon(x)=\mat I+\varepsilon \mat B(x),
\qquad \varepsilon\ge 0.
\]
Suponga adem\'as que
\[
\|\mat B\|_{L^\infty(\Omega)}\le M,
\qquad
M>0,
\qquad
0\le \varepsilon<\frac1M,
\]
y que $c\in L^\infty(\Omega)$, $\beta\in L^\infty(\Gamma_R)$, con $c,\beta\ge 0$. Suponga tambi\'en
\[
    f\in L^2(\Omega),
    \qquad
    r\in H^{-1/2}(\Gamma_R),
    \qquad
    g\in H^{1/2}(\Gamma_D).
\]

Considere el problema
\[
\begin{aligned}
-\dive\!\big(\mat A_\varepsilon \grad u\big)+cu &= f &&\text{en }\Omega,\\
\gamma_0u &= g &&\text{en }\Gamma_D,\\
(\mat A_\varepsilon\grad u)\cdot \vec n+\beta\,\gamma_0u &= r &&\text{en }\Gamma_R.
\end{aligned}
\]
Defina
\[
V=\{v\in H^1(\Omega):\gamma_0v=0\text{ en }\Gamma_D\},
\qquad
V_g=\{v\in H^1(\Omega):\gamma_0v=g\text{ en }\Gamma_D\}.
\]

\begin{enumerate}[label=(\alph*)]
    \item Tome un lifting $G\in H^1(\Omega)$ tal que $\gamma_0G=g$ en $\Gamma_D$, escriba $u=u_0+G$ con $u_0\in V$, y derive la formulaci\'on d\'ebil equivalente para $u_0$. Identifique con cuidado la forma bilineal $a_\varepsilon$ y la funcional lineal $\widetilde \ell_\varepsilon$.

    \item Demuestre que $a_\varepsilon$ es continua y coerciva en $V$, y concluya existencia y unicidad de $u_0$ por Lax--Milgram. Explique por qu\'e esto implica existencia y unicidad de $u\in V_g$.

    \item Suponga ahora que $u\in H^2(\Omega)$ y que $\mat A_\varepsilon\in W^{1,\infty}(\Omega)$. Partiendo de la formulaci\'on d\'ebil, recupere:
    \begin{itemize}
        \item la ecuaci\'on fuerte en $\Omega$,
        \item la condici\'on de Robin/conormal en $\Gamma_R$,
        \item y explique por qu\'e la condici\'on de Dirichlet en $\Gamma_D$ ya estaba incorporada desde el comienzo.
    \end{itemize}

    \item Discuta brevemente qu\'e puede fallar si $\Gamma_D=\varnothing$, $c=0$ y $\beta=0$. Relacione esto con el caso de Neumann puro.
\end{enumerate}

\end{enumerate}

\end{document}
